\chapter{Persistent index và quản lý cấu hình}
\label{ch:persist}

\section{Vì sao chọn}

Embedding 54 Node trên CPU mất khoảng 77 giây; dựng lại mỗi lần hỏi là không chấp nhận được. Ngược lại,
dùng lại chỉ mục cũ sau khi đã đổi mô hình embedding hay cách cắt Node thì cho kết quả sai mà không báo
lỗi. Cần lưu chỉ mục \emph{và} biết khi nào nó hết hợp lệ.

\section{Cơ chế}

Hàm \texttt{build\_or\_load\_index()} gọi \texttt{storage\_context.persist()} sau khi dựng chỉ mục; ở các lần
sau, nó tạo \texttt{StorageContext} từ thư mục đã lưu rồi gọi \texttt{load\_index\_from\_storage()}. Thư mục lưu
gồm docstore, vector store và index store dạng JSON.

\subsection{Fingerprint: chống dùng lại chỉ mục cũ}

Bản trước chỉ lưu tên mô hình embedding và sự tồn tại của \texttt{docstore.json}, nên sửa cách cắt
Node mà quên \texttt{--rebuild} sẽ âm thầm dùng chỉ mục cũ. Bản mới lưu \texttt{corpus.json} cạnh
chỉ mục, gồm: model embedding, thư mục corpus, \textbf{phiên bản loader cùng các ngưỡng cắt}, và
sha256 + tên hiển thị của từng file. Chỉ cần một trong các thứ đó đổi là chỉ mục tự build lại.

\begin{lstlisting}[language=Python, caption={Fingerprint của chỉ mục (\texttt{src/index\_builder.py}).}, label={lst:fingerprint}]
    return {
        "embed_model": cfg.embed_model_name,
        "corpus_dir": cfg.corpus_dir,
        "loader": {"phien_ban": LOADER_VERSION, "max_tu": MAX_TU, "min_tu_noi_dung": MIN_TU_NOI_DUNG},
        "files": files,
    }
\end{lstlisting}

Từ bản này, \texttt{persist\_dir} được suy ra thành \texttt{storage/<profile>/<đường dẫn corpus>}, nên hai profile có số chiều vector khác nhau không còn ghi đè lên nhau (bản báo cáo trước ghi nhận đây là việc ``nên làm''). Vì số chiều vector gắn với mô hình, đổi mô hình embedding vẫn phải build lại chỉ mục --- \texttt{corpus.json} cạnh chỉ mục sẽ tự phát hiện.


\section{Số liệu}

\begin{table}[H]
  \centering
  \small
  \begin{tabularx}{\textwidth}{X l}
    \toprule
    \textbf{Thao tác (profile \texttt{server}, CPU, 54 Node)} & \textbf{Thời gian} \\
    \midrule
    Dựng mới (embedding toàn bộ Node) & khoảng 77 giây \\
    Nạp lại từ đĩa (kể cả nạp mô hình embedding) & 1,5 giây (đo ngày 09/10) \\
    \bottomrule
  \end{tabularx}
  \caption{Thời gian dựng mới và nạp lại chỉ mục.}
\end{table}

\section{Dùng ở đâu trong đồ án}

Mọi điểm vào (\texttt{pipeline.py}, \texttt{app/app.py}, \texttt{eval/*.py}, notebook) đều gọi
\texttt{build\_or\_load\_index()}, nên không có đường nào dùng nhầm chỉ mục cũ.
